\begin{resumo}
O presente documento estabelece a modelagem formal e a especificação detalhada de Casos de Uso do Sistema Inteligente de Gestão de Eventos Adversos baseado na metodologia \textit{Global Trigger Tool} (SIGEA-GTT), desenvolvido no âmbito da Universidade Federal de Sergipe (UFS). Fundamentado nos preceitos da Engenharia de Software Orientada a Casos de Uso e nas diretrizes formais da especificação OMG UML 2.5, o trabalho delimita a fronteira do sistema e caracteriza com precisão os quatro atores operacionais da plataforma: o Administrador de TI (\texttt{ADMIN}), o Docente Orientador e Auditor Sênior (\texttt{PROFESSOR}), o Discente Auditor em Treinamento (\texttt{STUDENT}) e o Ator Temporal do Sistema (\texttt{SYSTEM\_TIMER}). São concebidos tanto a visão macro do sistema quanto quatro diagramas modulares em linguagem vetorial nativa TikZ, abrangendo segurança e governança de acesso com validação obrigatória institucional (\texttt{@academico.ufs.br}), parametrização canônica do catálogo IHI-GTT (6 módulos e 53 gatilhos) e escala NCC MERP de dano, gestão pedagógica de turmas acadêmicas, execução de auditoria clínica cronometrada em 20 minutos com prontuário simulado em formato estruturado semiestruturado JSONB, ferramentas da qualidade para análise de causa-raiz (Diagrama de Ishikawa 6M, Matriz GUT, Plano de Ação 5W2H e Ciclo PDCA), e avaliação docente formativa com consolidação de indicadores epidemiológicos. Um total de vinte e dois casos de uso é detalhado de forma exaustiva em conformidade com o padrão de Cockburn, contemplando fluxos principais, alternativos e de exceção, pré e pós-condições, regras de negócio invioláveis e restrições não funcionais. Por fim, são apresentadas matrizes completas de rastreabilidade bidirecional com os vinte e cinco requisitos funcionais (RF01 a RF25), atributos de qualidade da ISO/IEC 25010 e entidades físico-relacionais do banco de dados PostgreSQL 16 LTS, assegurando a integridade e a viabilidade da implementação do sistema.

\vspace{\onelineskip}
\noindent
\textbf{Palavras-chave}: Engenharia de Software. Casos de Uso. UML 2.5. Global Trigger Tool. Gestão de Eventos Adversos. SIGEA-GTT. DCOMP/UFS.
\end{resumo}
